\begin{abstract}
\noindent
Picking $K$ assets from a universe of $n$ to balance risk and return is a textbook problem that turns NP-hard the moment cardinality is enforced, and it has become a flagship test case for near-term quantum optimisation. We apply the Quantum Approximate Optimisation Algorithm \parencite{farhi2014qaoa} to the cardinality-constrained mean--variance portfolio selection problem of Markowitz, implemented from scratch as a pure-NumPy statevector simulator without external quantum frameworks. The objective is reformulated as a QUBO with budget penalty $A$, mapped to an Ising cost Hamiltonian, and minimised by a variational circuit alternating cost and transverse-field mixer unitaries, with COBYLA driving ten random multi-starts. The canonical instance fixes $n=16$ equities drawn from mag7, quantum, quantum-big, and anti baskets at $K=4$, $\lambda=2.0$, $A=0.5$, and depth $p=3$. Four sweeps probe depth $p$, problem size $n$, risk aversion $\lambda$, and penalty $A$, complemented by a full $p=1$ energy-landscape map (\cref{subsec:landscape-results}) and brute-force, greedy-Sharpe, rounded-Markowitz, and simulated-annealing baselines. QAOA reaches a scaled ratio $r\approx 0.94$, yet the unscaled relative gap stays of order $10^{3}$ and the top-1 success probability collapses from $\sim\!0.23$ at $n=4$ to $\sim\!10^{-4}$ at $n=16$. Classical baselines hit the optimum at machine precision roughly $10^{3}$--$10^{4}\times$ faster. No quantum advantage is claimed; the contribution is methodological.
\end{abstract}
